\section{Questão 2}

\subsection{Enunciado}

Considere o seguinte conjunto de dados:
\begin{verbatim}
data = {
    'Tamanho': [150, 120, 180, 130, 160],
    'Quartos': [3, 2, 4, 3, 4],
    'Preço': [200000, 180000, 250000, 190000, 220000],
    'Localização': [
        'Norte', 'Leste', 'Oeste', 'Sul', 'Centro'
    ]
}
\end{verbatim}
Aplique a técnica de seleção de features para selecionar as features mais relevantes para o negócio.

\subsection{Solução}

O alvo é \texttt{Preço}. As entradas são \texttt{Tamanho}, \texttt{Quartos} e \texttt{Localização}. \texttt{Localização} entra como one-hot. A seleção usa \texttt{SelectKBest} com \texttt{f\_regression} e \(k = 2\): a estatística \(F\) da regressão linear de cada feature com o preço. As duas de maior \(F\) são mantidas.

\subsection{Resultado}

Na amostra, \texttt{Tamanho} (\(F = 54{,}56\)) e \texttt{Quartos} (\(F = 10{,}62\)) apresentaram os maiores F-scores e foram selecionadas pelo \texttt{SelectKBest}. \texttt{Localização} foi codificada em indicadores one-hot e avaliada, mas não entrou entre as duas primeiras. A Tabela~\ref{tab:q2-rank} traz o ranking completo.

\begin{table}[H]
    \centering
    \caption{Ranking das features pela estatística \(F\) em relação ao preço.}
    \label{tab:q2-rank}
    \footnotesize
    \setlength{\tabcolsep}{6pt}
    \begin{tabular}{@{}lccc@{}}
        \toprule
        Feature & \(F\) & \(p\) & Selecionada \\
        \midrule
        Tamanho            & 54,56 & 0,005 & Sim \\
        Quartos            & 10,62 & 0,047 & Sim \\
        Localização\_Oeste &  7,56 & 0,071 & Não \\
        Localização\_Leste &  1,40 & 0,322 & Não \\
        Localização\_Sul   &  0,45 & 0,549 & Não \\
        Localização\_Centro&  0,19 & 0,695 & Não \\
        Localização\_Norte &  0,08 & 0,796 & Não \\
        \bottomrule
    \end{tabular}
\end{table}
